\subsection{Main interpretation}
We study anatomical concepts as explicit interfaces between image features and fine-grained grading. The position posterior and lesion map are more inspectable than an undifferentiated latent feature, and their predictions are available at deployment without privileged metadata. The controlled comparison shows that their joint use has a distinct and compelling empirical pattern: either concept alone degrades the primary risk, whereas the joint configuration recovers baseline-level performance and is best observed.

The pattern is well aligned with the structure of the task. The position concept supplies an acquisition-aware anatomical frame, because the liver's appearance and fibrosis-related cues depend on the ultrasound position. The lesion concept supplies local pathological content within that organ-level frame. Together they expose a complementary representation of where the scan is taken and what lesion evidence is present, allowing the grading pathway to combine spatial context with local pathology. This frame-and-content view gives the joint improvement a principled anatomical rationale and turns the result into an interpretable scientific story rather than a purely numerical ranking. The explicit bidirectionally detached interface makes this proposed mechanism concrete and ready for targeted concept interventions.

\subsection{Interpretability claim and next test}
The model is best described as an anatomically grounded, non-bottleneck concept pathway. It is ``intervention-ready'' because the predicted concept representations enter through an explicit interface. A stronger interpretability claim would require replacing predicted concepts with ground-truth, corrupted, shuffled, or clinician-corrected concepts and measuring the resulting change in grading. Such an experiment would test sufficiency, necessity, and correction utility directly. We do not label the current attention maps as explanations merely because they are visually plausible.

\subsection{Why a single ResNet-50 is useful here}
Using one ResNet-50 backbone keeps the comparison controlled and makes the concept pathway the variable of interest. It is sufficient for a proof of concept in this ultrasound task. It does not support a claim that the interface is architecture-independent. Multi-seed and cross-backbone studies are natural follow-ups, but they are not silently inferred from the present result.

\subsection{Clinical meaning and limitations}
The additional position and weak-lesion outputs may help a reader check acquisition protocol adherence and inspect image regions associated with the score. This is a potential workflow benefit, not a clinical utility study. The cohort is restricted to schistosomiasis-associated fibrosis, the test centers are represented in training, and labels are expert ultrasound assessments rather than histopathology. The weak boxes are coarse and do not define a segmentation reference. The formal comparison uses one seed and one backbone, and the small E--A difference is not a demonstrated improvement.

\subsection{Conclusion}
SynAP-Fib provides a controlled example of converting training-time anatomical supervision into image-derived concepts that a fine-grained grading model can consume and expose. Its current evidence supports a task-specific concept pathway with best-observed joint performance and modest overhead. It does not yet establish faithful explanations, causal concept use, or transfer beyond the studied disease and backbone. Direct concept interventions are the next decisive experiment.
